\scanpage{145}
% archive foliation: 144
classes d'isomorphie d'objets d'icelle est en corr. biunivoque canonique avec l'ens des applications holomorphes
\begin{equation}
\tau:S\longrightarrow\mathbb C\setminus\mathbb R.
\tag{5}
\end{equation}
On dit qu'une structure de Torelli est ``directe'' ou ``stricte'' si \deleted{$t$ est} pour tt $s\in S$, la base \deleted{de} $z_1,z_2$ de $\mathcal H_s\otimes_{\mathbb Z}\mathbb R$ \deleted{correspond à}\inserted{définit} l'\deleted{\unsure{structure}}orientation de cet espace correspondant à sa structure complexe, i.e. si (5) est un morphisme
\begin{equation}
\tau:S\longrightarrow\mathcal D^+=\{z\in\mathbb C\mid\Im z>0\}.
\tag{6}
\end{equation}
Donc $\mathcal D^+$ apparaît comme l'espace analytique modulaire pour la structure ``courbe elliptique relative avec str. de Torelli stricte sur espace analytique variable'', la famille modulaire étant donnée par
\begin{equation}
\unsure{\mathcal X|\mathfrak X}_{\mathcal D^+}=(\mathcal D^+\times\mathbb C)/\mathbb Z^2
\tag{7}
\end{equation}
où $\mathbb Z^2$ opère sur $\mathcal D^+\times\mathbb C$ par
\begin{equation}
\deleted{\unsure{u}}\ T_{m,n}(\tau,z)=(\tau,z+m+n\tau)
\qquad\tau\in\mathcal D^+,\ z\in\mathbb C
\tag{8}
\end{equation}
de sorte que l'on a
\begin{equation}
(\unsure{\mathcal X|\mathfrak X}_{\mathcal D^+})_\tau
\simeq\mathbb C/(\mathbb Z+\mathbb Z\cdot\tau)
\simeq\mathbb C^*/\substack{\text{sous-groupe}\\
\text{engendré}\\
\text{par }\deleted{\unsure{q}}\ \unsure{\kappa|k}=\exp(2i\pi\tau)\\
(\text{satisfaisant }|\unsure{\kappa|k}|<1)}
\tag{8}
\end{equation}
